%\needspace{4\baselineskip}
\color{uniblue}{\subsection{Защита от перегрузки токами высших гармоник}}
\color{black}

\begin{enumerate}[label=\arabic{section}.\arabic{subsection}.\arabic{enumi}, labelsep=4pt, leftmargin=0pt, itemindent=57pt, itemsep=0pt, parsep=0pt] % если что можно крутить itemsep


\item
Защита от перегрузки предназначена для предотвращения перегрузки БСК токами высших 
гармоник. Измерительные цепи защиты включены на токи фаз А, В, С со стороны высоковольтных 
вводов БСК. Данная функция защиты реагирует на величину полного тока (истинное 
среднеквадратичное значение), включающую основную гармонику и высокочастотные гармонические 
составляющие тока. 

\item
Защита содержит сигнальную и отключающие ступени. Каждая из ступеней имеет регулируемое 
значение по току срабатывания, задаваемое параметром \textbf{«Iсигн ЗПВГ»} и \textbf{«Iоткл ЗПВГ»} для сигнальной 
и отключающей ступени соответственно. Сигнальная ступень с задержкой по времени \textbf{«Тсигн ЗПВГ»} 
действует на сигнализацию, отключающая ступень с задержкой по времени \textbf{«Тоткл ЗПВГ»} действует на 
отключение выключателя БСК. 

\begin{figure}[H]
\centering
\includegraphics[width=1\textwidth,height=1\textheight,keepaspectratio]{\fbpath/261.0381 ЗПВГ/img/ZPVG.pdf}
\captionof{figure}{Функциональный блок ЗПВГ}\label{pict:ZPVG}
\end{figure}




\small
\setlength{\extrarowheight}{0.05cm} % высота ячеек в таблице
\begin{longtable}
{|>{\centering\arraybackslash}m{4.5cm}|>{\centering\arraybackslash}m{4.5cm}|>{\centering\arraybackslash}m{3cm}|>{\centering\arraybackslash}m{3.5cm}|}
\caption{Перечень уставочных параметров ЗПВГ\hfill\vspace{-0.5\baselineskip}}\label{ksn:tbl2} \\ 
\hhline{|-|-|-|-|}
\rowcolor{gray!30}
Наименование уставки & Диапазон & Шаг &  Значение по умолчанию \\ 
\hhline{|-|-|-|-|}
\endfirsthead

\caption*{\hspace{3pt}\emph{Продолжение таблицы \ref{ksn:tbl2}\hfill\vspace{-0.5\baselineskip}}} \\ % сделано по ГОСТ 2.105 п.6.8.7
\hhline{|-|-|-|-|}
\rowcolor{gray!30}
Наименование уставки & Диапазон & Шаг &  Значение по умолчанию \\  
\hhline{|-|-|-|-|}
\endhead
\hhline{|-|-|-|-|}

\endfoot
\endlastfoot

%%%%%%%%%%%%содержимое%%%%%%%%%%%%%%

\end{longtable}
\setlength{\extrarowheight}{0.0cm} % возврат параметра высота ячеек в таблице
\normalsize







\end{enumerate} 